\documentclass[12pt, titlepage]{article}

%Page Size
\usepackage{geometry}
\geometry{a4paper}

\usepackage{amsmath, mathtools}
\usepackage{amsfonts}
\usepackage{amssymb}
\usepackage{graphicx}
\usepackage{colortbl}
\usepackage{xr}
\usepackage{hyperref}
\usepackage{longtable}
\usepackage{xfrac}
\usepackage{tabularx}
\usepackage{float}
\usepackage{siunitx}
\usepackage{booktabs}
\usepackage{caption}
\usepackage{pdflscape}
\usepackage{afterpage}

\hypersetup{
    colorlinks,
    citecolor=blue,
    filecolor=black,
    linkcolor=red,
    urlcolor=blue
}
%\usepackage[round]{natbib}

\input{../Comments}
%% Common Parts

\newcommand{\progname}{Cut Edge Weight: Data Separability Measurement} % PUT YOUR PROGRAM NAME HERE
\newcommand{\authname}{Lesley Wheat} % AUTHOR NAMES                  

\usepackage{hyperref}
    \hypersetup{colorlinks=true, linkcolor=blue, citecolor=blue, filecolor=blue,
                urlcolor=blue, unicode=false}
    \urlstyle{same}
                                


\begin{document}

\title{Project Title: System Verification and Validation Plan for \progname{}} 
\author{\authname}
\date{\today}
	
\maketitle

\pagenumbering{roman}

\section{Revision History}

\begin{tabularx}{\textwidth}{p{3cm}p{2cm}X}
\toprule {\bf Date} & {\bf Version} & {\bf Notes}\\
\midrule
2023-02-10 & 0.0 & Creation\\
2023-02-14 & 1.0 & Version 1.0\\
\bottomrule
\end{tabularx}

\newpage

\tableofcontents

\newpage

\listoftables
\mycomment{Remove this section if it isn't needed}

%\listoffigures
\mycomment{Remove this section if it isn't needed}

\newpage

\section{Symbols, Abbreviations and Acronyms}

\renewcommand{\arraystretch}{1.2}
\begin{tabular}{l l} 
  \toprule		
  \textbf{symbol} & \textbf{description}\\
  \midrule 
  T & Test\\
  CEWS & Cut Edge Weight Statistic\\
  $J_N$ & Normalized Multiclass Cut Edge Weight Measurement\\
  $D$ & Input Data Matrix\\
  $C$ & Input Class Matrix\\
  FR & Functional Requirement \\
  NFR & Non Functional Requirement \\
  \bottomrule
\end{tabular}\\

\mycomment{symbols, abbreviations or acronyms --- you can simply reference the SRS
  \citep{SRS} tables, if appropriate}

\mycomment{Remove this section if it isn't needed}

\newpage

\pagenumbering{arabic}

This document outlines the verification and validation plan for testing the CEWS System.

\section{General Information}

\subsection{Summary}

This document contains the verification and validation plan for CEWS System, a program for calculating the CEWS measurement for a dataset as shown in \cite{zighed_statistical_2005}.
CEWS System takes a set of data points $D$ and a list of class labels $C$ as inputs and computes the Normalized Multiclass Cut Edge Weight Measurement $J_N$.

\subsection{Objectives}

The objective of this document is to build confidence in the software’s correctness.
This document provides the test plan used to verify and validate the final product.
Therefore, it should aim to encompass all validation and verification components required for testing.

\subsection{Relevant Documentation}

Relevant documentation includes the system requirements specification found here: \href{https://github.com/LesleyWheat/CAS741_CEWS/tree/main/docs/SRS}{GitHub Link}.

\section{Plan}

\mycomment{Introduce this section.   You can provide a roadmap of the sections to
  come.}

\subsection{Verification and Validation Team} \label{sec:VnV_team}

\mycomment{Your teammates.  Maybe your supervisor.
  You shoud do more than list names.  You should say what each person's role is
  for the project's verification.  A table is a good way to summarize this information.}

The verification and validation team includes the domain expert Joachim de Fourestier, the verification and validation reviewer Mina Mahdipour, and Dr. Spencer Smith.

\subsection{SRS Verification Plan}

The SRS will be reviewed by the SRS review team, consisting of the domain expert Joachim de Fourestier and the SRS reviewer, Sam Joseph Crawford, and Dr. Spencer Smith.

\subsection{Design Verification Plan}

The design document MIS will be reviewed by the MIS review team, consisting of the domain expert Joachim de Fourestier and the MIS reviewer, Deesha Vipulkumar Patel, and Dr. Spencer Smith.

\subsection{Verification and Validation Plan Verification Plan}

This document will be reviewed by the VnV team members in Section \ref{sec:VnV_team}.

\mycomment{The verification and validation plan is an artifact that should also be verified.}

\mycomment{The review will include reviews by your classmates}

\mycomment{Create a checklists?}

\subsection{Implementation Verification Plan}

Individual tests are outlined in sections 5.1 through 5.2. 
The tests will either be performed manually or with someone reviewing the code.

\mycomment{You should at least point to the tests listed in this document and the unit
  testing plan.}

\mycomment{In this section you would also give any details of any plans for static verification of
  the implementation.  Potential techniques include code walkthroughs, code
  inspection, static analyzers, etc.}

\subsection{Automated Testing and Verification Tools}
The VnV proccess will use no tools.

\mycomment{What tools are you using for automated testing.  Likely a unit testing
  framework and maybe a profiling tool, like ValGrind.  Other possible tools
  include a static analyzer, make, continuous integration tools, test coverage
  tools, etc.  Explain your plans for summarizing code coverage metrics.
  Linters are another important class of tools.  For the programming language
  you select, you should look at the available linters.  There may also be tools
  that verify that coding standards have been respected, like flake9 for
  Python.}

\mycomment{If you have already done this in the development plan, you can point to
that document.}

\mycomment{The details of this section will likely evolve as you get closer to the
  implementation.}

\subsection{Software Validation Plan}

To test the output of the software, results from the paper \cite{zighed_statistical_2005} will be used as a pseudo-oracle.
The following benchmark datasets from the UCI Machine Learning Repository will be used: Breast Cancer, BUPA liver, Glass Ident, Haberman, Image seg, Ionosphere, Iris (Bezdek), Iris plants, Musk "Clean1", Pima Indians, Waveform, Wine recognition and Yeast.

\mycomment{If there is any external data that can be used for validation, you should
  point to it here.  If there are no plans for validation, you should state that
  here.}

\mycomment{You might want to use review sessions with the stakeholder to check that
the requirements document captures the right requirements.  Maybe task based
inspection?}

\mycomment{This section might reference back to the SRS verification section.}

%---------------------------------------------------------------------------------
\section{System Test Description}
	
\subsection{Tests for Functional Requirements}

\mycomment{Subsets of the tests may be in related, so this section is divided into
  different areas.  If there are no identifiable subsets for the tests, this
  level of document structure can be removed.}

\mycomment{Include a blurb here to explain why the subsections below
  cover the requirements.  References to the SRS would be good here.}

\subsubsection{Area of Testing 1: FR1, FR4}

This area of testing will address inputs to the program.

Functional requirement 1 states: CEWS System will read the input data. Also functional requirement 4 states: The CEWS System will generate an output.

\mycomment{It would be nice to have a blurb here to explain why the subsections below
  cover the requirements.  References to the SRS would be good here.  If a section
  covers tests for input constraints, you should reference the data constraints
  table in the SRS.}
		
\paragraph{Tests for FR1}

\begin{enumerate}

\item{T1 - Normal Input\\}

Control: Manual
					
Initial State: Pending Input
					
Input: $D$ and $C$ have the same number of samples and features, $C$ contains at least two classes.
					
Output: No error generated, program proceeds to next state.

Test Case Derivation: FR1, FR4
					
How test will be performed: Manual
					
\item{T2 - Data Mismatch\\}

Control: Manual
					
Initial State: 
					
Input: $D$ and $C$ have different number of samples or $D$ has a different number of features.
					
Output: Generate error about incorrect data format.

Test Case Derivation: FR1, FR4

How test will be performed: Manual

\item{T3 - Single Class\\}

Control: Manual
					
Initial State: 
					
Input: $C$ contains only one class
					
Output: Generate error about lack of classes.

Test Case Derivation: FR1, FR4

How test will be performed: Manual

\item{T4 - Empty Data\\}

Control: Manual
					
Initial State: Pending Input
					
Input: $C$ or $D$ is null.
					
Output: Generate warning about empty data object.

Test Case Derivation: FR1, FR4

How test will be performed: Manual

\item{T5 - Value out of range\\}

Control: Manual
					
Initial State: Pending Input
					
Input: $C$ contains a value larger than $W_{\max}$ or smaller than $W_{\min}$. $D$ contains a value larger than $F_{\max}$ or smaller than $F_{\min}$. (Table \ref{Tabel:constants})
					
Output: Generate warning about data value out of range.

Test Case Derivation: FR1, FR4

How test will be performed: Manual

\end{enumerate}

\subsubsection{Area of Testing 2: FR2, FR3, FR4}
This area of testing addresses the result of the program.
Functional Requirement 2 states: CEWS System will preform the calculation.

Functional Requirement 3 states: The CEWS System should generate an error if the calculation result exceeds the constraints.

Functional Requirement 4 states: The CEWS System will generate an output.

\begin{enumerate}

\item{T6 - UCI Dataset Tests\\}

Control: Manual
					
Initial State: Pending Input
					
Input: UCI Machine Learning Repository datasets: Breast Cancer, BUPA liver, Glass Ident, Haberman, Image seg, Ionosphere, Iris (Bezdek), Iris plants, Musk "Clean1", Pima Indians, Waveform, Wine recognition and Yeast.
					
Output: Normalized CEWS Measurement $J_N$.

Test Case Derivation: FR2, FR3, FR4
					
How test will be performed:
\begin{enumerate}
\item Format a dataset
\item Input the reformatted dataset to the system
\item Check the output value against the results in Table \ref{Table:benchmarks}.
\end{enumerate}

\begin{table}[H]
\centering
\begin{tabular}{|c|c|} \hline
Dataset & Expected Result \\ \hline \hline
\iref{Breast Cancer} & 0.008\\ \hline
\rref{BUPA liver} & 0.401 \\ \hline
\rref{Glass Ident.} & 0.356 \\ \hline
\rref{Haberman}   & 0.331 \\ \hline
\rref{Image seg.} & 0.224 \\ \hline
\rref{Ionosphere} & 0.137\\ \hline
\rref{Iris (Bezdek)} & 0.090 \\ \hline
\rref{Iris Plants} & 0.087 \\ \hline
\rref{Musk "Clean1"} & 0.167 \\ \hline
\rref{Pima Indians} & 0.310 \\ \hline
\rref{Waveform} & 0.255 \\ \hline
\rref{Wine recognition} & 0.093\\ \hline
\rref{Yeast} & 0.524 \\ \hline
\end{tabular}
\caption{CEWS Result on Benchmark Datasets from \cite{zighed_statistical_2005}}
\label{Table:benchmarks}
\end{table}

\end{enumerate}


\subsection{Tests for Nonfunctional Requirements}

\mycomment{The nonfunctional requirements for accuracy will likely just reference the
  appropriate functional tests from above.  The test cases should mention
  reporting the relative error for these tests.  Not all projects will
  necessarily have nonfunctional requirements related to accuracy}

\mycomment{Tests related to usability could include conducting a usability test and
  survey.  The survey will be in the Appendix.}

\mycomment{Static tests, review, inspections, and walkthroughs, will not follow the
format for the tests given below.}

\subsubsection{Area of Testing 1: NFR1}

Non Functional Requirement 1 states: The accuracy of the computed solutions should meet the level needed for data analysis.
The absolute error should be less than 0.01.

		
%\paragraph{Title for Test}

\begin{enumerate}

\item{T7: Accuracy\\}

Type: Manual
					
Initial State: Pending input
					
Input/Condition: UCI Datasets listed in Table \ref{Table:benchmarks}.
					
Output/Result: Less than 0.01 difference in results from Table \ref{Table:benchmarks}.
					
How test will be performed: Manual

\end{enumerate}

\subsubsection{Area of Testing 2: NFR2}

Non Functional Requirement 2 states: A program should be able to call CEWS System by calling a function of the form CEWS($D$, $C$) and CEWS System should return the output as described in FR4.

\begin{enumerate}

\item{T8: Usability\\}

Type: Manual
					
Initial State: Off
					
Input/Condition: Simulated dataset
					
Output/Result: $J_N$
					
How test will be performed: Manually, the user will attempt to call the program's function.

\end{enumerate}

\subsubsection{Area of Testing 3: NFR3}

Non Functional Requirement 3 states: The program code should be readable and clear.

\begin{enumerate}

\item{T9: Maintainability\\}

Type: Manual
					
Initial State: None
					
Input/Condition: None
					
Output/Result: None
					
How test will be performed: The author will review the code.

\end{enumerate}

\subsubsection{Area of Testing 4: NFR4}

Non Functional Requirement 4 states: The software should run on the Windows 11 operating system

\begin{enumerate}

\item{T9: Maintainability\\}

Type: Manual
					
Initial State: Off
					
Input/Condition: Simulated dataset
					
Output/Result: $J_N$
					
How test will be performed: Manually on a computer running Windows 11.

\end{enumerate}

\subsection{Traceability Between Test Cases and Requirements}

\mycomment{Provide a table that shows which test cases are supporting which
  requirements.}

%\afterpage{
%\begin{landscape}
\begin{table}[h!]
\centering
\begin{tabular}{|c|c|c|c|c|c|c|c|c|c|}
\hline
	& T1 & T2 & T3 & T4 & T5 & T6 & T7 & T8 & T9\\
\hline
FR1  &X&X&X&X&X& & & &  \\ \hline
FR2  & & & & &X& & & &  \\ \hline
FR3  & & & & &X& & & & \\ \hline
FR4  & &X&X&X&X& & & &  \\ \hline
NFR1 & & & & & &X& & &  \\ \hline
NFR2 & & & & & & &X& &  \\ \hline
NFR3 & & & & & & & &X&  \\ \hline
NFR4 & & & & & & & & &X  \\ \hline
\end{tabular}
\caption{Which test cases are supporting which requirements}
\label{Table:test_support}
\end{table}
%\end{landscape}
%}


\section{Unit Test Description}

This section is intentionally left blank until the MIS is completed.

\mycomment{
\mycomment{Reference your MIS (detailed design document) and explain your overall
  philosophy for test case selection.}  
\mycomment{This section should not be filled in until after the MIS (detailed design
  document) has been completed.}

\subsection{Unit Testing Scope}

\mycomment{What modules are outside of the scope.  If there are modules that are
  developed by someone else, then you would say here if you aren't planning on
  verifying them.  There may also be modules that are part of your software, but
  have a lower priority for verification than others.  If this is the case,
  explain your rationale for the ranking of module importance.}

\subsection{Tests for Functional Requirements}

\mycomment{Most of the verification will be through automated unit testing.  If
  appropriate specific modules can be verified by a non-testing based
  technique.  That can also be documented in this section.}

\subsubsection{Module 1}

\mycomment{Include a blurb here to explain why the subsections below cover the module.
  References to the MIS would be good.  You will want tests from a black box
  perspective and from a white box perspective.  Explain to the reader how the
  tests were selected.}

\begin{enumerate}

\item{test-id1\\}

Type: \mycomment{Functional, Dynamic, Manual, Automatic, Static etc. Most will
  be automatic}
					
Initial State: 
					
Input: 
					
Output: \mycomment{The expected result for the given inputs}

Test Case Derivation: \mycomment{Justify the expected value given in the Output field}

How test will be performed: 
					
\item{test-id2\\}

Type: \mycomment{Functional, Dynamic, Manual, Automatic, Static etc. Most will
  be automatic}
					
Initial State: 
					
Input: 
					
Output: \mycomment{The expected result for the given inputs}

Test Case Derivation: \mycomment{Justify the expected value given in the Output field}

How test will be performed: 

\item{...\\}
    
\end{enumerate}

\subsubsection{Module 2}

...

\subsection{Tests for Nonfunctional Requirements}

\mycomment{If there is a module that needs to be independently assessed for
  performance, those test cases can go here.  In some projects, planning for
  nonfunctional tests of units will not be that relevant.}

\mycomment{These tests may involve collecting performance data from previously
  mentioned functional tests.}

\subsubsection{Module ?}
		
\begin{enumerate}

\item{test-id1\\}

Type: \mycomment{Functional, Dynamic, Manual, Automatic, Static etc. Most will
  be automatic}
					
Initial State: 
					
Input/Condition: 
					
Output/Result: 
					
How test will be performed: 
					
\item{test-id2\\}

Type: Functional, Dynamic, Manual, Static etc.
					
Initial State: 
					
Input: 
					
Output: 
					
How test will be performed: 

\end{enumerate}

\subsubsection{Module ?}

...

\subsection{Traceability Between Test Cases and Modules}
}
\mycomment{Provide evidence that all of the modules have been considered.}

\newpage
\phantomsection
\addcontentsline{toc}{section}{References}
\bibliographystyle{IEEEtran}
\bibliography{refs/theoryrefs}

\newpage

\section{Appendix}

This is where you can place additional information.

\subsection{Symbolic Parameters}

The definition of the test cases will call for SYMBOLIC\_CONSTANTS.
Their values are defined in this section for easy maintenance.

\begin{table}[H]
\caption{Table of SYMBOLIC\_CONSTANTS} \label{Tabel:constants}
\renewcommand{\arraystretch}{1.2}
\noindent \begin{longtable*}{l l} 
  \toprule
  \textbf{Var} & \textbf{Value} \\
  \midrule 
  $F_{\min}$ & $-(2 -2^{-23})2^{127}$\\
  $F_{\max}$ & $(2-2^{-23})2^{127}$\\
    $W_{\min}$ & 0\\
  $W_{\max}$ &  $2^{8} - 1$\\
  $J_{\max}$ & 1\\
  $J_{\min}$ & 0\\
  \bottomrule
\end{longtable*}
\end{table}

%\subsection{Usability Survey Questions?}

\mycomment{This is a section that would be appropriate for some projects.}

\mycomment{
\newpage{}
\section*{Appendix --- Reflection}

The information in this section will be used to evaluate the team members on the
graduate attribute of Lifelong Learning.  Please answer the following questions:

\newpage{}
\section*{Appendix --- Reflection}

The information in this section will be used to evaluate the team members on the
graduate attribute of Lifelong Learning.  Please answer the following questions:

\begin{enumerate}
  \item What knowledge and skills will the team collectively need to acquire to
  successfully complete the verification and validation of your project?
  Examples of possible knowledge and skills include dynamic testing knowledge,
  static testing knowledge, specific tool usage etc.  You should look to
  identify at least one item for each team member.
  \item For each of the knowledge areas and skills identified in the previous
  question, what are at least two approaches to acquiring the knowledge or
  mastering the skill?  Of the identified approaches, which will each team
  member pursue, and why did they make this choice?
\end{enumerate}
}
\end{document}